% Kết luận (bản rút gọn). Rút từ sections/conclusion.tex; mọi con số đã xuất hiện ở các mục trước, trừ các
% tỉ số khuyến nghị (nguồn ở comment cuối đoạn khuyến nghị của sections/conclusion.tex).
Nhóm cài đặt ba spinlock, một sense-reversing barrier và ba họ concurrent counter, kiểm tra chúng bằng
stress test, ThreadSanitizer và mutation testing, rồi đo chúng trên ba máy. RQ2 tới RQ5 dựa trên
\PCname{apple-m5}. RQ1 và RQ6 có số liệu của cả ba máy.

RQ1. Không bài kiểm tra nào tìm thấy lỗi trên ba máy, và bộ kiểm tra bắt được cả 14 biến thể cố ý làm
hỏng. Lỗi thuật toán lộ ra dưới stress test, còn lỗi memory order chủ yếu do ThreadSanitizer bắt.
Tính đúng của các cài đặt dựa trên kết quả này cùng lập luận memory order ở mục Cài đặt.

RQ2. TTAS mở rộng tốt nhất trong ba spinlock tự viết (2,48 lần cận dưới ở 20 thread), còn lock thư viện
không sụp vì thread chờ ngủ trong kernel. MCS gãy mạnh tại số core và sụp khi vượt số core, vì nó trao
lock cho đúng một thread kế tiếp, kể cả khi thread đó không có core.

RQ3. Backoff giúp TAS nhiều nhất (ở 20 thread, TAS với sleep backoff chỉ tốn 0,31 lần thời gian của
TAS), nhưng TTAS với sleep backoff chậm hơn TTAS 6,1--6,8 lần khi vượt số core với $\mathrm{OW} = 3$.
Lợi ích của backoff đến từ việc giữ lock ở một chủ, nên lock nhanh nhất là lock kém công bằng nhất.

RQ4. Với barrier spin, một episode tốn thêm khoảng 53~ns cho mỗi thread trong cụm P-core, và số episode
mỗi giây giảm khoảng 2\,000 lần khi số thread vượt số core. Khi đó, chờ bằng \texttt{sched\_yield}
nhanh nhất.

RQ5. Từ 4 thread, fetch-and-add nhanh hơn vòng CAS 2,0 tới 2,7 lần, vì số lần CAS hỏng trên mỗi lần
tăng tăng từ 0,09 ở 2 thread lên 1,59 ở 8 thread. Khi không có việc giữa hai lần tăng, CAS có backoff
nhanh hơn fetch-and-add từ 2 thread. Với sleep backoff, các thread tăng lần lượt (độ công bằng xấp xỉ
$1/p$); với spin và yield backoff, chỉ một phần các lần tăng dồn về một thread. Còn
\PCname{lock_std_mutex} và \PCname{lock_pthread_mutex} ngang fetch-and-add ở 4 thread và nhanh hơn từ 6
thread.

RQ6. Một phần, vì thứ hạng chính giữ trên cả ba máy nhưng giữa hai máy Apple, điểm gãy dời theo số
P-core và điểm sụp của MCS dời theo số core. Trên \PCname{LAP1184}, \texttt{seq\_cst} không tốn thêm
gì, barrier spin chậm 17 lần ngay ở 12 thread, và \texttt{sleep\_for} với yêu cầu 100~ns ngủ trung vị
khoảng 62~\textmu s.

Phép đo khoảng trống trao lock cho thấy điểm gãy và điểm sụp của MCS đến từ các lần trao lock phải chờ
scheduler, không từ lưu lượng cache coherence. Điểm gãy tại số P-core phù hợp với giả thuyết hai loại
core, nhưng chưa có phép đo riêng cho cơ chế này. Cơ chế khiến counter dùng mutex thư viện thắng
fetch-and-add, và cơ chế khiến \texttt{seq\_cst} làm CAS chậm khoảng 2 lần trên \PCname{apple-m5}, cũng
mới là giả thuyết.

Khi số thread có thể vượt số core, nhóm khuyến nghị dùng lock thư viện. Trên macOS,
\PCname{os_unfair_lock} nhanh nhất trong các lock thư viện, chậm 1,6--1,7 lần cận dưới khi không có
việc ngoài critical section. Tránh TTAS với sleep backoff khi có việc ngoài critical section. Khi tranh
chấp cao, chỉ dùng MCS khi số thread không vượt số P-core. Trên máy chạy trực tiếp, dùng barrier spin
khi số thread không vượt số core, và barrier có \texttt{sched\_yield} khi vượt. Dùng fetch-and-add cho
counter khi có việc giữa hai lần tăng.

Việc chưa làm gồm phiên bản phân tán bằng MPI-3 RMA (phần 4 của đề), counter chia theo thread (mỗi
thread một counter riêng, cộng lại khi đọc), một workload có ghi dữ liệu dùng chung trong critical
section, và một máy x86 chạy Linux trực tiếp để tách ảnh hưởng của máy ảo.
